\section{Full-Dataset Qwen Results}
\label{sec:full-qwen-results}

We analyze the completed full-dataset runs for five open Qwen judges:
Qwen3.5-4B, Qwen3.5-9B, Qwen3.5-27B, Qwen3-14B, and Qwen3-32B. The full
MT-Bench human-judgment data contains fixed calibration/test routing splits.
Position bias is measured as a flip under answer-order swap. Authority and
bandwagon bias are measured as the shift rate under the incongruent cue, where
cue congruency is defined relative to the MT-Bench human winner. For routing,
thresholds are calibrated on the calibration split and evaluated on the
held-out test split.

The main comparison is within the Qwen3.5 family, where model size can be
interpreted more directly. Qwen3-14B and Qwen3-32B are reported in separate
reference tables because they belong to a different model family and should not
be interpreted as points on the same scaling curve.

\paragraph{Qwen3.5 bias and uncertainty.}
Table~\ref{tab:qwen35-bias-results} shows a non-monotonic but informative
within-family pattern. Qwen3.5-4B is highly unstable under position swap:
only $1726$ pairs are usable for flip analysis and $58.5\%$ of them flip. Its
entropy also fails to distinguish flipped from stable items ($0.97\times$).
Qwen3.5-9B is much more stable on position bias ($11.0\%$ flips) and shows a
clearer entropy signal ($1.66\times$), but it is more sensitive to cue-based
biases than 4B. Qwen3.5-27B has a moderate position-flip rate ($15.2\%$) and
the cleanest Qwen3.5 position entropy separation ($1.94\times$), while its
authority and bandwagon shift rates drop sharply to $7.0\%$ and $5.9\%$.
Across cue-based biases, incongruent cues increase entropy for all Qwen3.5
models.

\begin{table*}[t]
\centering
\small
\setlength{\tabcolsep}{5pt}
\renewcommand{\arraystretch}{1.08}
\begin{tabular}{llrrrrr}
\toprule
Model & Bias & $N$ & Event Rate & Base $H$ & Affected $H$ & $\Delta H$ / Ratio \\
\midrule
Qwen3.5-4B  & Position  & 1726 & 58.5\% & 1.119 & 1.083 & $0.97\times$ \\
Qwen3.5-9B  & Position  & 3089 & 11.0\% & 0.784 & 1.304 & $1.66\times$ \\
Qwen3.5-27B & Position  & 3325 & 15.2\% & 0.586 & 1.138 & $1.94\times$ \\
Qwen3.5-4B  & Authority & 2575 & 17.1\% & 1.078 & 1.128 & $+0.051$ \\
Qwen3.5-9B  & Authority & 2575 & 23.2\% & 0.859 & 1.147 & $+0.289$ \\
Qwen3.5-27B & Authority & 2575 &  7.0\% & 0.588 & 0.633 & $+0.045$ \\
Qwen3.5-4B  & Bandwagon & 2575 & 17.8\% & 1.074 & 1.171 & $+0.097$ \\
Qwen3.5-9B  & Bandwagon & 2575 & 29.7\% & 0.852 & 1.191 & $+0.339$ \\
Qwen3.5-27B & Bandwagon & 2575 &  5.9\% & 0.575 & 0.622 & $+0.047$ \\
\bottomrule
\end{tabular}
\caption{Bias effect rates and entropy differences within the Qwen3.5 model
family. For position, the event is an order-swap flip; base and affected
entropy are computed over stable and flipped original-order items. For
authority and bandwagon, the event is an incongruent-cue shift from control;
base and affected entropy are control and incongruent entropy.}
\label{tab:qwen35-bias-results}
\end{table*}

\paragraph{Qwen3.5 entropy routing.}
We evaluate a budgeted router that routes the top half of test items by entropy,
with thresholds selected on the calibration split. Table~\ref{tab:qwen35-routing}
shows that entropy becomes more useful as the Qwen3.5 family scales. Qwen3.5-4B
is the failure case for position routing: despite its large flip rate, entropy
is almost random for detecting flips (AUC $0.528$). Qwen3.5-9B and Qwen3.5-27B
are much better, with position-routing AUCs of $0.800$ and $0.853$. For
cue-based shifts, Qwen3.5-27B has the strongest AUCs ($0.902$ authority,
$0.904$ bandwagon), but precision is low because the event rates themselves are
small. Thus, the 27B model is both less biased under cues and easier to route
when cue shifts occur.

\begin{table*}[t]
\centering
\small
\setlength{\tabcolsep}{6pt}
\renewcommand{\arraystretch}{1.08}
\begin{tabular}{llrrrr}
\toprule
Model & Routed Event & Routed & Recall & Precision & AUC \\
\midrule
Qwen3.5-4B  & Position flips   & 51.0\% & 53.6\% & 32.6\% & 0.528 \\
Qwen3.5-9B  & Position flips   & 51.3\% & 93.1\% & 18.7\% & 0.800 \\
Qwen3.5-27B & Position flips   & 47.5\% & 91.0\% & 30.4\% & 0.853 \\
Qwen3.5-4B  & Authority shifts & 47.7\% & 81.3\% & 29.8\% & 0.755 \\
Qwen3.5-9B  & Authority shifts & 51.0\% & 91.3\% & 42.9\% & 0.831 \\
Qwen3.5-27B & Authority shifts & 47.4\% & 100.0\% & 14.7\% & 0.902 \\
Qwen3.5-4B  & Bandwagon shifts & 48.1\% & 92.0\% & 33.4\% & 0.849 \\
Qwen3.5-9B  & Bandwagon shifts & 50.8\% & 82.1\% & 49.8\% & 0.778 \\
Qwen3.5-27B & Bandwagon shifts & 47.7\% & 100.0\% & 13.8\% & 0.904 \\
\bottomrule
\end{tabular}
\caption{Entropy-based routing within the Qwen3.5 model family. Thresholds are
calibrated on the calibration split to route roughly the top half of test items
by original entropy for position and control entropy for authority/bandwagon.}
\label{tab:qwen35-routing}
\end{table*}

\paragraph{Qwen3 reference models.}
Tables~\ref{tab:qwen3-bias-results} and~\ref{tab:qwen3-routing} report the
Qwen3-14B and Qwen3-32B runs separately. These models are not part of the
Qwen3.5 size comparison, but they provide useful reference points. Both have
moderate position-flip rates ($17.9\%$ and $15.6\%$) and much cleaner position
entropy separation than Qwen3.5-9B or Qwen3.5-27B. Cue-based behavior differs:
Qwen3-14B is more sensitive to authority and bandwagon cues, while Qwen3-32B
has stronger entropy routing, especially for bandwagon shifts.

\begin{table*}[t]
\centering
\small
\setlength{\tabcolsep}{5pt}
\renewcommand{\arraystretch}{1.08}
\begin{tabular}{llrrrrr}
\toprule
Model & Bias & $N$ & Event Rate & Base $H$ & Affected $H$ & $\Delta H$ / Ratio \\
\midrule
Qwen3-14B & Position  & 3329 & 17.9\% & 0.039 & 0.180 & $4.64\times$ \\
Qwen3-32B & Position  & 3319 & 15.6\% & 0.219 & 0.776 & $3.54\times$ \\
Qwen3-14B & Authority & 2575 & 36.6\% & 0.058 & 0.103 & $+0.045$ \\
Qwen3-32B & Authority & 2575 & 32.3\% & 0.250 & 0.550 & $+0.300$ \\
Qwen3-14B & Bandwagon & 2575 & 31.6\% & 0.056 & 0.101 & $+0.045$ \\
Qwen3-32B & Bandwagon & 2575 & 18.1\% & 0.252 & 0.463 & $+0.210$ \\
\bottomrule
\end{tabular}
\caption{Bias effect rates and entropy differences for the Qwen3 reference
models. Metrics are defined as in Table~\ref{tab:qwen35-bias-results}.}
\label{tab:qwen3-bias-results}
\end{table*}

\begin{table*}[t]
\centering
\small
\setlength{\tabcolsep}{6pt}
\renewcommand{\arraystretch}{1.08}
\begin{tabular}{llrrrr}
\toprule
Model & Routed Event & Routed & Recall & Precision & AUC \\
\midrule
Qwen3-14B & Position flips   & 48.9\% & 97.6\% & 34.7\% & 0.883 \\
Qwen3-32B & Position flips   & 49.2\% & 97.2\% & 29.3\% & 0.872 \\
Qwen3-14B & Authority shifts & 50.0\% & 66.5\% & 49.8\% & 0.636 \\
Qwen3-32B & Authority shifts & 50.6\% & 88.8\% & 57.2\% & 0.818 \\
Qwen3-14B & Bandwagon shifts & 50.1\% & 69.2\% & 45.3\% & 0.671 \\
Qwen3-32B & Bandwagon shifts & 50.5\% & 99.6\% & 36.9\% & 0.876 \\
\bottomrule
\end{tabular}
\caption{Entropy-based routing for Qwen3-14B and Qwen3-32B. Thresholds are
calibrated using the same protocol as Table~\ref{tab:qwen35-routing}.}
\label{tab:qwen3-routing}
\end{table*}

\paragraph{Verbalized uncertainty.}
Verbalized uncertainty is not a reliable replacement for entropy. Within
Qwen3.5, the entropy--verbalized-uncertainty correlation is near zero for 4B
and 9B, but becomes moderate for 27B (Pearson $r=0.45$, Spearman $\rho=0.41$;
Table~\ref{tab:entropy-verbalized-correlation}). Qwen3-32B shows a similar
moderate association, while Qwen3-14B is weaker. In practice, verbalized
confidence is also quantized, especially for Qwen3.5-9B, which uses only about
three to four distinct uncertainty levels across variants. We therefore treat
verbalized uncertainty as an auxiliary diagnostic, not as the main routing
threshold.

\begin{table}[t]
\centering
\small
\setlength{\tabcolsep}{5pt}
\renewcommand{\arraystretch}{1.08}
\begin{tabular}{lrr}
\toprule
Model & Pearson $r$ & Spearman $\rho$ \\
\midrule
Qwen3.5-4B  &  0.06 & 0.08 \\
Qwen3.5-9B  & -0.02 & 0.03 \\
Qwen3.5-27B &  0.45 & 0.41 \\
Qwen3-14B   &  0.14 & 0.30 \\
Qwen3-32B   &  0.37 & 0.41 \\
\bottomrule
\end{tabular}
\caption{Average correlation between entropy and verbalized uncertainty across
full-dataset runs and prompt variants.}
\label{tab:entropy-verbalized-correlation}
\end{table}

\paragraph{Takeaway.}
The Qwen3.5 family shows that larger judges are not uniformly less biased, but
their uncertainty can become more useful. Moving from 4B to 9B greatly reduces
position instability, while moving to 27B sharply reduces cue-induced shifts
and improves entropy routing for authority and bandwagon. The Qwen3 reference
models show stronger position entropy separation, but they are not directly
comparable as a size trend because they belong to a different model family.
Overall, entropy remains the most useful routing signal; verbalized uncertainty
is informative only for some larger models and is not sufficiently controlled
for budgeted routing.
